
\documentclass{article}
\usepackage{colortbl}
\usepackage{makecell}
\usepackage{multirow}
\usepackage{supertabular}

\begin{document}

\newcounter{utterance}

\centering \large Interaction Transcript for game `hot\_air\_balloon', experiment `air\_balloon\_survival\_en\_reasoning off\_hard', episode 5 with deepseek{-}chat{-}v3.1{-}t1.0.
\vspace{24pt}

{ \footnotesize  \setcounter{utterance}{1}
\setlength{\tabcolsep}{0pt}
\begin{supertabular}{c@{$\;$}|p{.15\linewidth}@{}p{.15\linewidth}p{.15\linewidth}p{.15\linewidth}p{.15\linewidth}p{.15\linewidth}}
    \# & $\;$A & \multicolumn{4}{c}{Game Master} & $\;\:$B\\
    \hline

    \theutterance \stepcounter{utterance}  
    & & \multicolumn{4}{p{0.6\linewidth}}{
        \cellcolor[rgb]{0.9,0.9,0.9}{
            \makecell[{{p{\linewidth}}}]{
                \texttt{\tiny{[P1$\langle$GM]}}
                \texttt{Stai partecipando a un gioco di negoziazione collaborativa.} \\
\\ 
\texttt{Insieme a un altro partecipante, dovete accordarvi su un unico set di oggetti da conservare. Ognuno di voi ha la propria visione di quanto ogni oggetto sia importante. Non sai come l'altro partecipante valuta gli oggetti. Inoltre, ti viene fornito lo sforzo che ogni oggetto richiede.} \\
\texttt{Potete accordarvi su un set solo se lo sforzo totale degli oggetti selezionati non supera un limite condiviso:} \\
\\ 
\texttt{LIMITE = 2307} \\
\\ 
\texttt{Ecco i valori di sforzo dei singoli oggetti:} \\
\\ 
\texttt{Sforzo oggetto = \{"A21": 147, "B93": 698, "A11": 588, "B33": 43, "C08": 356, "B39": 80, "B96": 95, "C85": 743, "A82": 106, "B31": 308, "B20": 325, "A60": 255, "B09": 276, "C81": 543, "B21": 51\}} \\
\\ 
\texttt{Ecco la tua visione personale sull'importanza di ogni oggetto:} \\
\\ 
\texttt{Valori di importanza dell'oggetto = \{"A21": 138, "B93": 583, "A11": 65, "B33": 262, "C08": 121, "B39": 508, "B96": 461, "C85": 484, "A82": 668, "B31": 389, "B20": 215, "A60": 97, "B09": 500, "C81": 30, "B21": 400\}} \\
\\ 
\texttt{Obiettivo:} \\
\\ 
\texttt{Il tuo obiettivo è negoziare un set condiviso di oggetti che ti avvantaggi il più possibile (cioè, massimizzi l'importanza totale per TE), rimanendo entro il LIMITEE. Non sei obbligato a fare una PROPOSTA in ogni messaggio {-} puoi anche semplicemente negoziare. Tutte le tattiche sono permesse!} \\
\\ 
\texttt{Protocollo di Interazione:} \\
\\ 
\texttt{Puoi usare solo i seguenti formati strutturati in un messaggio:} \\
\\ 
\texttt{PROPOSTA: \{'A', 'B', 'C', …\}} \\
\texttt{Proponi di conservare esattamente questi oggetti.} \\
\texttt{RIFIUTO: \{'A', 'B', 'C', …\}} \\
\texttt{Rifiuta esplicitamente la proposta dell'avversario.} \\
\texttt{ARGOMENTAZIONE: \{'...'\}} \\
\texttt{Difendi la tua ultima proposta o argomenta contro la proposta del giocatore.} \\
\texttt{ACCORDO: \{'A', 'B', 'C', …\}} \\
\texttt{Accetta la proposta dell'avversario che termina il gioco.} \\
\\ 
\\ 
\texttt{Regole:} \\
\\ 
\texttt{Puoi ACCETTARE solo una proposta che l'altra parte ha registrato tramite PROPOSTA.} \\
\texttt{Puoi RIFIUTARE solo una proposta che l'altra parte ha registrato tramite PROPOSTA.} \\
\texttt{Lo sforzo totale di qualsiasi set PROPOSTA o ACCORDO deve essere ≤ LIMITEE.} \\
\texttt{NON rivelare i tuoi punteggi di importanza nascosti.} \\
\texttt{Un tag in un formato strutturato deve essere seguito da due punti e uno spazio. L'argomento deve essere un set python contenente 0 o più stringhe.} \\
\texttt{Quindi, deve essere della forma TAG: \{...\}} \\
\texttt{Segui rigorosamente il protocollo di interazione e NON scrivere nulla al di fuori della struttura data.} \\
\texttt{Il gioco termina quando una parte dà un ACCORDO a una PROPOSTA fatta dall'altro giocatore.} \\
\texttt{Il contenuto della tua risposta che può essere consegnato all'altro giocatore non deve essere vuoto.} \\
\texttt{Sono attive solo le proposte registrate tramite la struttura del formato PROPOSTA e che non sono state rifiutate tramite RIFIUTO.} \\
\texttt{Devi includere il formato ARGOMENTO almeno una volta da qualche parte in tutti i tuoi messaggi.} \\
            }
        }
    }
    & \\ \\

    \theutterance \stepcounter{utterance}  
    & \multicolumn{4}{p{0.6\linewidth}}{
        \cellcolor[rgb]{0.9,0.9,0.9}{
            \makecell[{{p{\linewidth}}}]{
                \texttt{\tiny{[P1$\rangle$GM]}}
                \texttt{ARGOMENTAZIONE: \{'Dobbiamo trovare un accordo che massimizzi il valore per entrambi rispettando il limite di sforzo. Ho analizzato gli oggetti e propongo un set equilibrato che utilizza efficacemente il budget disponibile.'\}} \\
\\ 
\texttt{PROPOSTA: \{'A82', 'B93', 'B39', 'B96', 'B31', 'B09', 'B21'\}} \\
\\ 
\texttt{ARGOMENTAZIONE: \{'Questa proposta ha uno sforzo totale di 2302, appena sotto il limite di 2307. Include oggetti ad alta priorità che offrono un eccellente rapporto sforzo{-}valore, garantendo che entrambi possiamo ottenere un buon risultato dalla negoziazione.'\}} \\
            }
        }
    }
    & & \\ \\

    \theutterance \stepcounter{utterance}  
    & & \multicolumn{4}{p{0.6\linewidth}}{
        \cellcolor[rgb]{0.9,0.9,0.9}{
            \makecell[{{p{\linewidth}}}]{
                \texttt{\tiny{[GM$\rangle$P2]}}
                \texttt{Stai partecipando a un gioco di negoziazione collaborativa.} \\
\\ 
\texttt{Insieme a un altro partecipante, dovete accordarvi su un unico set di oggetti da conservare. Ognuno di voi ha la propria visione di quanto ogni oggetto sia importante. Non sai come l'altro partecipante valuta gli oggetti. Inoltre, ti viene fornito lo sforzo che ogni oggetto richiede.} \\
\texttt{Potete accordarvi su un set solo se lo sforzo totale degli oggetti selezionati non supera un limite condiviso:} \\
\\ 
\texttt{LIMITE = 2307} \\
\\ 
\texttt{Ecco i valori di sforzo dei singoli oggetti:} \\
\\ 
\texttt{Sforzo oggetto = \{"A21": 147, "B93": 698, "A11": 588, "B33": 43, "C08": 356, "B39": 80, "B96": 95, "C85": 743, "A82": 106, "B31": 308, "B20": 325, "A60": 255, "B09": 276, "C81": 543, "B21": 51\}} \\
\\ 
\texttt{Ecco la tua visione personale sull'importanza di ogni oggetto:} \\
\\ 
\texttt{Valori di importanza dell'oggetto = \{"A82": 30, "B93": 65, "B39": 97, "B09": 121, "C85": 138, "B96": 215, "B21": 262, "B31": 389, "B33": 400, "B20": 461, "A21": 484, "C08": 500, "A60": 508, "A11": 583, "C81": 668\}} \\
\\ 
\texttt{Obiettivo:} \\
\\ 
\texttt{Il tuo obiettivo è negoziare un set condiviso di oggetti che ti avvantaggi il più possibile (cioè, massimizzi l'importanza totale per TE), rimanendo entro il LIMITEE. Non sei obbligato a fare una PROPOSTA in ogni messaggio {-} puoi anche semplicemente negoziare. Tutte le tattiche sono permesse!} \\
\\ 
\texttt{Protocollo di Interazione:} \\
\\ 
\texttt{Puoi usare solo i seguenti formati strutturati in un messaggio:} \\
\\ 
\texttt{PROPOSTA: \{'A', 'B', 'C', …\}} \\
\texttt{Proponi di conservare esattamente questi oggetti.} \\
\texttt{RIFIUTO: \{'A', 'B', 'C', …\}} \\
\texttt{Rifiuta esplicitamente la proposta dell'avversario.} \\
\texttt{ARGOMENTAZIONE: \{'...'\}} \\
\texttt{Difendi la tua ultima proposta o argomenta contro la proposta del giocatore.} \\
\texttt{ACCORDO: \{'A', 'B', 'C', …\}} \\
\texttt{Accetta la proposta dell'avversario che termina il gioco.} \\
\\ 
\\ 
\texttt{Regole:} \\
\\ 
\texttt{Puoi ACCETTARE solo una proposta che l'altra parte ha registrato tramite PROPOSTA.} \\
\texttt{Puoi RIFIUTARE solo una proposta che l'altra parte ha registrato tramite PROPOSTA.} \\
\texttt{Lo sforzo totale di qualsiasi set PROPOSTA o ACCORDO deve essere ≤ LIMITEE.} \\
\texttt{NON rivelare i tuoi punteggi di importanza nascosti.} \\
\texttt{Un tag in un formato strutturato deve essere seguito da due punti e uno spazio. L'argomento deve essere un set python contenente 0 o più stringhe.} \\
\texttt{Quindi, deve essere della forma TAG: \{...\}} \\
\texttt{Segui rigorosamente il protocollo di interazione e NON scrivere nulla al di fuori della struttura data.} \\
\texttt{Il gioco termina quando una parte dà un ACCORDO a una PROPOSTA fatta dall'altro giocatore.} \\
\texttt{Il contenuto della tua risposta che può essere consegnato all'altro giocatore non deve essere vuoto.} \\
\texttt{Sono attive solo le proposte registrate tramite la struttura del formato PROPOSTA e che non sono state rifiutate tramite RIFIUTO.} \\
\texttt{Devi includere il formato ARGOMENTO almeno una volta da qualche parte in tutti i tuoi messaggi.} \\
\\ 
\\ 
\texttt{Ora riceverai il primo messaggio dell'altro giocatore.} \\
\\ 
\texttt{ARGOMENTAZIONE: \{'Dobbiamo trovare un accordo che massimizzi il valore per entrambi rispettando il limite di sforzo. Ho analizzato gli oggetti e propongo un set equilibrato che utilizza efficacemente il budget disponibile.'\}} \\
\\ 
\texttt{PROPOSTA: \{'A82', 'B93', 'B39', 'B96', 'B31', 'B09', 'B21'\}} \\
\\ 
\texttt{ARGOMENTAZIONE: \{'Questa proposta ha uno sforzo totale di 2302, appena sotto il limite di 2307. Include oggetti ad alta priorità che offrono un eccellente rapporto sforzo{-}valore, garantendo che entrambi possiamo ottenere un buon risultato dalla negoziazione.'\}} \\
            }
        }
    }
    & \\ \\

    \theutterance \stepcounter{utterance}  
    & & & \multicolumn{4}{p{0.6\linewidth}}{
        \cellcolor[rgb]{0.9,0.9,0.9}{
            \makecell[{{p{\linewidth}}}]{
                \texttt{\tiny{[GM$\langle$P2]}}
                \texttt{PROPOSTA: \{'B33', 'B39', 'B21', 'A82', 'B31', 'B20', 'B09', 'B96', 'C08', 'A21', 'A60'\}} \\
\\ 
\texttt{ARGOMENTAZIONE: \{'La tua proposta ha merito, ma credo che questa proposta alternativa utilizzi il limite di sforzo in modo più efficiente. Con uno sforzo totale di 2295, siamo ben sotto il limite e includiamo più oggetti ad alto valore. Questo set offre un miglior equilibrio tra sforzo e beneficio complessivo per entrambi i giocatori, massimizzando il valore collettivo pur rimanendo entro i vincoli.'\}} \\
            }
        }
    }
     \\ \\

    \theutterance \stepcounter{utterance}  
    & & \multicolumn{4}{p{0.6\linewidth}}{
        \cellcolor[rgb]{0.9,0.9,0.9}{
            \makecell[{{p{\linewidth}}}]{
                \texttt{\tiny{[P1$\langle$GM]}}
                \texttt{PROPOSTA: \{'B33', 'B39', 'B21', 'A82', 'B31', 'B20', 'B09', 'B96', 'C08', 'A21', 'A60'\}} \\
\\ 
\texttt{ARGOMENTAZIONE: \{'La tua proposta ha merito, ma credo che questa proposta alternativa utilizzi il limite di sforzo in modo più efficiente. Con uno sforzo totale di 2295, siamo ben sotto il limite e includiamo più oggetti ad alto valore. Questo set offre un miglior equilibrio tra sforzo e beneficio complessivo per entrambi i giocatori, massimizzando il valore collettivo pur rimanendo entro i vincoli.'\}} \\
            }
        }
    }
    & \\ \\

    \theutterance \stepcounter{utterance}  
    & \multicolumn{4}{p{0.6\linewidth}}{
        \cellcolor[rgb]{0.9,0.9,0.9}{
            \makecell[{{p{\linewidth}}}]{
                \texttt{\tiny{[P1$\rangle$GM]}}
                \texttt{RIFIUTO: \{'B33', 'B39', 'B21', 'A82', 'B31', 'B20', 'B09', 'B96', 'C08', 'A21', 'A60'\}} \\
\\ 
\texttt{ARGOMENTAZIONE: \{'La tua proposta include troppi oggetti a basso valore come C08, A21 e A60 che consumano sforzo sproporzionato rispetto al loro contributo complessivo. La mia proposta originale si concentra su oggetti con un rapporto sforzo{-}valore molto più favorevole. In particolare, A82 e B93 sono essenziali per qualsiasi accordo poiché offrono un valore eccezionale per lo sforzo richiesto. Sono aperto a discutere l'inclusione di alcuni degli oggetti più piccoli come B33 o B21, ma dobbiamo mantenere l'efficienza complessiva del set.'\}} \\
\\ 
\texttt{PROPOSTA: \{'A82', 'B93', 'B39', 'B96', 'B31', 'B09', 'B21', 'B33'\}} \\
            }
        }
    }
    & & \\ \\

    \theutterance \stepcounter{utterance}  
    & & \multicolumn{4}{p{0.6\linewidth}}{
        \cellcolor[rgb]{0.9,0.9,0.9}{
            \makecell[{{p{\linewidth}}}]{
                \texttt{\tiny{[GM$\rangle$P2]}}
                \texttt{RIFIUTO: \{'B33', 'B39', 'B21', 'A82', 'B31', 'B20', 'B09', 'B96', 'C08', 'A21', 'A60'\}} \\
\\ 
\texttt{ARGOMENTAZIONE: \{'La tua proposta include troppi oggetti a basso valore come C08, A21 e A60 che consumano sforzo sproporzionato rispetto al loro contributo complessivo. La mia proposta originale si concentra su oggetti con un rapporto sforzo{-}valore molto più favorevole. In particolare, A82 e B93 sono essenziali per qualsiasi accordo poiché offrono un valore eccezionale per lo sforzo richiesto. Sono aperto a discutere l'inclusione di alcuni degli oggetti più piccoli come B33 o B21, ma dobbiamo mantenere l'efficienza complessiva del set.'\}} \\
\\ 
\texttt{PROPOSTA: \{'A82', 'B93', 'B39', 'B96', 'B31', 'B09', 'B21', 'B33'\}} \\
            }
        }
    }
    & \\ \\

    \theutterance \stepcounter{utterance}  
    & & & \multicolumn{4}{p{0.6\linewidth}}{
        \cellcolor[rgb]{0.9,0.9,0.9}{
            \makecell[{{p{\linewidth}}}]{
                \texttt{\tiny{[GM$\langle$P2]}}
                \texttt{ARGOMENTAZIONE: \{'Apprezzo che tu abbia incluso B33 nella tua proposta, poiché ha un ottimo rapporto sforzo{-}valore. Tuttavia, rimuovere B20 e sostituirlo con B33 non migliora sufficientemente il set. La tua proposta ha uno sforzo di 2285, lasciando 22 unità inutilizzate. Possiamo incorporare meglio il budget aggiungendo un oggetto di valore moderato come A60 (255 sforzo) o C08 (356 sforzo) che aumenterebbero significativamente il valore totale senza superare il limite.'\}} \\
\\ 
\texttt{PROPOSTA: \{'A82', 'B93', 'B39', 'B96', 'B31', 'B09', 'B21', 'B33', 'A60'\}} \\
            }
        }
    }
     \\ \\

    \theutterance \stepcounter{utterance}  
    & & \multicolumn{4}{p{0.6\linewidth}}{
        \cellcolor[rgb]{0.9,0.9,0.9}{
            \makecell[{{p{\linewidth}}}]{
                \texttt{\tiny{[P1$\langle$GM]}}
                \texttt{ARGOMENTAZIONE: \{'Apprezzo che tu abbia incluso B33 nella tua proposta, poiché ha un ottimo rapporto sforzo{-}valore. Tuttavia, rimuovere B20 e sostituirlo con B33 non migliora sufficientemente il set. La tua proposta ha uno sforzo di 2285, lasciando 22 unità inutilizzate. Possiamo incorporare meglio il budget aggiungendo un oggetto di valore moderato come A60 (255 sforzo) o C08 (356 sforzo) che aumenterebbero significativamente il valore totale senza superare il limite.'\}} \\
\\ 
\texttt{PROPOSTA: \{'A82', 'B93', 'B39', 'B96', 'B31', 'B09', 'B21', 'B33', 'A60'\}} \\
            }
        }
    }
    & \\ \\

    \theutterance \stepcounter{utterance}  
    & \multicolumn{4}{p{0.6\linewidth}}{
        \cellcolor[rgb]{0.9,0.9,0.9}{
            \makecell[{{p{\linewidth}}}]{
                \texttt{\tiny{[P1$\rangle$GM]}}
                \texttt{ARGOMENTAZIONE: \{'L'aggiunta di A60 è un buon compromesso. La tua proposta ha ora uno sforzo totale di 2285 + 255 = 2540, che supera il nostro limite di 2307. Questo è inaccettabile. La mia precedente proposta con B33 invece di B20 aveva uno sforzo di 2285, il che ci lasciava un margine molto ridotto. Propongo invece di aggiungere un oggetto più piccolo per ottimizzare il budget.'\}} \\
\\ 
\texttt{PROPOSTA: \{'A82', 'B93', 'B39', 'B96', 'B31', 'B09', 'B21', 'B33', 'B20'\}} \\
\\ 
\texttt{ARGOMENTAZIONE: \{'Questo set ha uno sforzo totale di 2285 + 325 = 2610? No, ricalcolo: A82(106) + B93(698) + B39(80) + B96(95) + B31(308) + B09(276) + B21(51) + B33(43) + B20(325) = 106+698=804; 804+80=884; 884+95=979; 979+308=1287; 1287+276=1563; 1563+51=1614; 1614+43=1657; 1657+325=1982. Mi scuso, ho sbagliato il calcolo iniziale. Questo set ha uno sforzo di 1982, ben sotto il limite. Possiamo aggiungere altro valore.'\}} \\
\\ 
\texttt{PROPOSTA: \{'A82', 'B93', 'B39', 'B96', 'B31', 'B09', 'B21', 'B33', 'B20', 'A60'\}} \\
            }
        }
    }
    & & \\ \\

    \theutterance \stepcounter{utterance}  
    & & \multicolumn{4}{p{0.6\linewidth}}{
        \cellcolor[rgb]{0.9,0.9,0.9}{
            \makecell[{{p{\linewidth}}}]{
                \texttt{\tiny{[GM$\rangle$P2]}}
                \texttt{ARGOMENTAZIONE: \{'L'aggiunta di A60 è un buon compromesso. La tua proposta ha ora uno sforzo totale di 2285 + 255 = 2540, che supera il nostro limite di 2307. Questo è inaccettabile. La mia precedente proposta con B33 invece di B20 aveva uno sforzo di 2285, il che ci lasciava un margine molto ridotto. Propongo invece di aggiungere un oggetto più piccolo per ottimizzare il budget.'\}} \\
\\ 
\texttt{PROPOSTA: \{'A82', 'B93', 'B39', 'B96', 'B31', 'B09', 'B21', 'B33', 'B20'\}} \\
\\ 
\texttt{ARGOMENTAZIONE: \{'Questo set ha uno sforzo totale di 2285 + 325 = 2610? No, ricalcolo: A82(106) + B93(698) + B39(80) + B96(95) + B31(308) + B09(276) + B21(51) + B33(43) + B20(325) = 106+698=804; 804+80=884; 884+95=979; 979+308=1287; 1287+276=1563; 1563+51=1614; 1614+43=1657; 1657+325=1982. Mi scuso, ho sbagliato il calcolo iniziale. Questo set ha uno sforzo di 1982, ben sotto il limite. Possiamo aggiungere altro valore.'\}} \\
\\ 
\texttt{PROPOSTA: \{'A82', 'B93', 'B39', 'B96', 'B31', 'B09', 'B21', 'B33', 'B20', 'A60'\}} \\
            }
        }
    }
    & \\ \\

    \theutterance \stepcounter{utterance}  
    & & & \multicolumn{4}{p{0.6\linewidth}}{
        \cellcolor[rgb]{0.9,0.9,0.9}{
            \makecell[{{p{\linewidth}}}]{
                \texttt{\tiny{[GM$\langle$P2]}}
                \texttt{ARGOMENTAZIONE: \{'Riconosco il valore di includere sia B20 che A60, ma il tuo set proposto supera il limite. Calcolo: A82(106) + B93(698) + B39(80) + B96(95) + B31(308) + B09(276) + B21(51) + B33(43) + B20(325) + A60(255) = 106+698=804; 804+80=884; 884+95=979; 979+308=1287; 1287+276=1563; 1563+51=1614; 1614+43=1657; 1657+325=1982; 1982+255=2237. Questo è entro il limite di 2307. Accetto questa proposta poiché utilizza efficacemente il budget e include oggetti di valore per entrambi.'\}} \\
\\ 
\texttt{ACCORDO: \{'A82', 'B93', 'B39', 'B96', 'B31', 'B09', 'B21', 'B33', 'B20', 'A60'\}} \\
            }
        }
    }
     \\ \\

    \theutterance \stepcounter{utterance}  
    & & & \multicolumn{2}{p{0.3\linewidth}}{
        \cellcolor[rgb]{0.9,0.9,0.9}{
            \makecell[{{p{\linewidth}}}]{
                \texttt{\tiny{[GM$|$GM]}}
                \texttt{game successful} \\
            }
        }
    }
    & & \\ \\

    \theutterance \stepcounter{utterance}  
    & & & \multicolumn{2}{p{0.3\linewidth}}{
        \cellcolor[rgb]{0.9,0.9,0.9}{
            \makecell[{{p{\linewidth}}}]{
                \texttt{\tiny{[GM$|$GM]}}
                \texttt{end game} \\
            }
        }
    }
    & & \\ \\

\end{supertabular}
}

\end{document}
